\begin{center}
\fbox{\begin{minipage}{0.95\linewidth}
\small
\textbf{Key findings}
\begin{itemize}\setlength\itemsep{1pt}\setlength\parskip{0pt}
\item \textbf{``MaxSim'' is not one function, and quantization makes the difference large.} Libraries differ in normalization, epsilon, clamp and padding. On unit-length float vectors these rules nearly agree. On int4 tokens they diverge: NumKong's rule (contract C) keeps the top document of \sweet's rule (contract S) on only \nm{ce.beir.S-vs-C.top1}\% of BEIR calls.
\item \textbf{Two new defects and one known defect in public libraries.} Reference checks found invalid scores from NumKong on Intel CPUs with AMX and a padding clamp in maxsim-cpu's AArch64 build. They also confirmed a known maxsim-cpu defect on x86.
\item \textbf{NumKong's approximation failed when candidates scored close together.} In our tests, NumKong's rank correlation with the exact order fell from \nm{nkcheck.random.rho.min}--\nm{nkcheck.random.rho.max} to \nm{nkcheck.clustered.rho.min}--\nm{nkcheck.clustered.rho.max} when candidate scores lay close together.
\item \textbf{int4 storage can cost quality.} Keeping the length of decoded int4 tokens lost up to \nm{q.intel.nfcorpus.answersmall.relabs}\% nDCG@10 for a model with tight score margins; for ColBERTv2 the pre-registered test found no material loss.
\item \textbf{Speed depends on the condition, and the kernel is a small part of a search.} With warm calls on five CPUs, the int4 kernel of sweet-search was faster than every frozen exact float32 rival on captured agent and code-search calls. Rivals were faster with one thread, after a busy gap, with many callers in one process and on some large calls. In sweet-search's own search path, the kernel took a median of \nm{pipe.maxsim.med}\% of an agent search call.
\end{itemize}
\end{minipage}}
\end{center}
